% ===================================================================
\section{The Networked Individual \DPP\ Model}
% ===================================================================

\subsection{From Monolithic Supply Chain Records to Linked Passports}

A fundamental design choice in \DPP\ architecture is whether to represent a
product's entire supply chain as a single monolithic record or as a network of
linked individual passports. The current flat-file approach implicitly favours
the monolithic model: a single \DPP\ document attempts to encode all upstream
component data, material origins, and process information within a single
payload. This approach has several structural limitations. First, \textbf{data
staleness}: component data embedded at assembly time becomes stale if the
component's \DPP\ is updated later (e.g., a corrected hazard classification or a
supplier change). Second, \textbf{redundancy and inconsistency}: the same
component data may be duplicated across multiple parent product \DPPs, leading
to inconsistency when updates occur. Third, \textbf{size and complexity}: for
products with deep, multi-tiered supply chains (e.g., a battery electric vehicle
with thousands of components), a single monolithic \DPP\ becomes unwieldy.
Fourth, \textbf{modularity breakdown}: as a representative of \JTC~24 noted at
the \DPPforEU~2026 conference~\cite{dpp4eu2026_day1}, once a modular data model
reaches hundreds or thousands of data points, dependencies between modules
become extremely hard to track --- ``\textit{you don't really know what is going
to happen to the whole ecosystem}'' if a single referenced data point is
modified.

The Networked Individual \DPP\ model proposes an alternative: each product,
component, and material batch receives its own \DPP, and supply chain
relationships are represented as \textbf{semantic links} between these
individual passports. A parent product's \DPP\ does not embed its components'
data; it simply references their \DPP\ \URIs. This approach mirrors the
structure already recognised in the standards community. As a presenter from the
University of Pavia described at \DPPforEU~2026~\cite{dpp4eu2026_day1},
information units --- documents, top-level documents, sections, and parts of
sections --- are ``\textit{all represented uniformly as information units which
can cite one another. An information unit can have one parent at most. This
captures the hierarchical structure of standards and information units can be
derived from one or more parent or one or more source information units}''. This
description --- hierarchical structure, mutual citation, derivation from parent
units --- is structurally identical to the semantic linking model proposed here
for \DPPs.

\subsection{Core Semantic Linking Properties}

The Networked Individual \DPP\ model is built on three core semantic properties,
expressed in \RDF/\RDFS\ vocabulary consistent with the architecture's \RDF\
foundation (Section~3). The use of \texttt{rdf:Property} with
\texttt{rdfs:domain} and \texttt{rdfs:range} declarations --- rather than \OWL's
\texttt{owl:ObjectProperty} and \texttt{owl:DatatypeProperty} --- reflects the
deliberate choice of \RDF\ as the simplest, most generic representation for
\DPP\ data.

\subsubsection{\texttt{dpp:hasPrecursor} (The Assembly Link)}

\begin{itemize}
\item \textbf{Type}: \texttt{rdf:Property}
\item \textbf{Domain}: \texttt{dpp:ProductPassport}
\item \textbf{Range}: \texttt{dpp:ProductPassport}
\item \textbf{Purpose}: Links a parent product \DPP\ to the \DPPs\ of its
constituent components. It allows a query engine to crawl backwards through the
production network recursively.
\item \textbf{Cardinality}: One parent product may have many precursors; each
precursor has exactly one parent product.
\end{itemize}

\subsubsection{\texttt{dpp:derivedFrom} (The Transformation Link)}

\begin{itemize}
\item \textbf{Type}: \texttt{rdf:Property}
\item \textbf{Domain}: \texttt{dpp:ProductPassport}
\item \textbf{Range}: \texttt{dpp:ProductPassport}
\item \textbf{Purpose}: Used when a raw material is transformed (e.g., textile
yarn into a fabric roll, or raw chemical powders into a battery anode). It
establishes the semantic lineage without requiring the parent passport to
understand the internal chemistry of the input.
\item \textbf{Cardinality}: Many materials may derive from one or more source
materials.
\end{itemize}

\subsubsection{\texttt{dpp:massFraction} (The Quantification Attribute)}

\begin{itemize}
\item \textbf{Type}: \texttt{rdf:Property}
\item \textbf{Range}: \texttt{xsd:decimal}
\item \textbf{Purpose}: Declares exactly how much of the linked component's
weight or mass constitutes the final product. Associated with a blank node or
reified relation linking a precursor \DPP\ to its parent.
\item \textbf{Datatype}: \texttt{xsd:decimal} (e.g., 0.42 for 42\% mass
fraction).
\end{itemize}


\emph{Note: While the \DPP\ model is expressed in \RDF\ / \RDFS\ for maximum
simplicity, the properties are compatible with OWL and can be refined to
\texttt{owl:ObjectProperty} or \texttt{owl:DatatypeProperty} if \OWL\ reasoning
is required for specific applications.}

\subsection{\RDF\ / \JSONLD\ Representation}

When serialised, the parent passport remains clean, lightweight, and completely
generic. It simply maps its assembly topology. The critical insight is that the
parent \DPP\ contains \textbf{no component data} --- only links. The actual
composition, hazard, and provenance data for each component lives in its own
\DPP, stored behind the federated portal. When a query engine resolves
\texttt{dpp:hasPrecursor}, it follows the link \textit{within the federation of
national portals}, retrieving the component's data in real time. This design
directly addresses the modularity problem raised at the conference:
``\textit{once the size is limited to a few dozens, you can still fully
understand the overall data model. Once it gets to hundreds or thousands of data
points, that is impossible}''~\cite{dpp4eu2026_day1}. By distributing data
across linked individual \DPPs\ rather than accumulating it in a single
monolithic record, the model keeps each data unit manageable and independently
maintainable.

\begin{lstlisting}[language=jsonld]
{
  "@context": "https://w3id.org/dpp/v1",
  "@id": "urn:dpp:product:batch-2026-001",
  "@type": "dpp:ProductPassport",
  "dpp:productName": "Lithium-Ion Battery Pack Model X",
  "dpp:hasPrecursor": [
    {
      "@id": "urn:dpp:component:cathode-batch-789",
      "dpp:massFraction": {
        "@value": 0.42,
        "@type": "xsd:decimal"
      }
    },
    {
      "@id": "urn:dpp:component:anode-batch-456",
      "dpp:massFraction": {
        "@value": 0.28,
        "@type": "xsd:decimal"
      }
    },
    {
      "@id": "urn:dpp:component:electrolyte-batch-123",
      "dpp:massFraction": {
        "@value": 0.15,
        "@type": "xsd:decimal"
      }
    }
  ],
  "dpp:derivedFrom": {
    "@id": "urn:dpp:material:lithium-carbonate-batch-001"
  }
}
\end{lstlisting}

\subsection{Each \DPP\ as a Knowledge Graph (Named Graph)}

In the portal-backed architecture, each individual \DPP\ is not merely a
\JSONLD\ document but a \textbf{knowledge graph} stored as an \RDF\ Named Graph
within a national portal's triple-store. A Named Graph is an \RDF\ graph identified by
an \IRI, allowing the portal to manage each \DPP\ as a distinct, addressable
unit within a larger graph database. This design has several implications.

First, \textbf{independent lifecycle management}: each \DPP\ can be updated,
versioned, or deprecated independently of other \DPPs. When a component
manufacturer corrects a hazard classification, only that component's \DPP\ is
updated; all parent product \DPPs\ that reference it via
\texttt{dpp:hasPrecursor} automatically reflect the correction upon resolution,
without requiring synchronised updates across the supply chain.

Second, \textbf{\SPARQL\ queryability}: because each \DPP\ is a named graph
within a national portal's triple-store, each portal can execute cross-\DPP\
\SPARQL\ queries --- for example, traversing an entire supply chain from final
product to raw materials in a single query, or identifying all products
containing a specific substance of concern across all \DPPs\ in its own
triple-store and, on demand, across the whole federation of national portals.

Third, \textbf{access control granularity}: as discussed in Section~5, the
portal can enforce \ODRL\ access policies at the level of individual named
graphs or sub-graphs, providing branch-level access control that is not possible
when data is scattered across distributed \REO\ endpoints. The use of persistent
\IRIs\ as graph identifiers, as emphasised in the author's honorary lecture at
\DPPforEU~2026 where ``\textit{on the lowest level is the \IRI --- the
international resource identifier}''~\cite{dpp4eu2026_day2}, ensures that each
\DPP\ remains uniquely addressable and resolvable indefinitely.

\subsection{\IRI\ as the Language-Neutral Foundation}

Before a \DPP\ can be represented as a knowledge graph, its constituent entities
--- products, components, materials, substances, processes, organisations ---
must be unambiguously identified. The architecture adopts the International
Resource Identifier (\IRI) as the lowest level of identification, as emphasised
in the honorary lecture: ``\textit{on the lowest level is the \IRI --- the
international resource identifier}.'' The use of \IRIs\ resolves a fundamental
challenge in the European \DPP\ ecosystem: the multiplicity of natural languages
across member states. As Figure~\ref{fig:iri-multilingual} illustrates, the same
concept (``exit'') expressed in multiple natural languages
% --- 出口, Ausgang, Chūkǒu, sortie, exit, salida, uscita, thon sgo, keluar,
% kijárat, dalje ---
that
map to a single \IRI. This principle, applied throughout the \DPP\ architecture,
enables cross-lingual, cross-jurisdictional interoperability without requiring a
single natural language to dominate the system.

\begin{figure}[h]
\centering
\includegraphics[width=0.9\linewidth]{%
figures/2025-10-16__DigiPass_DPP_handling-IRI.pdf}
\caption{The \IRI\ as a language-neutral identifier. The same concept (``exit'')
expressed in multiple natural languages maps to a single International Resource
Identifier (IRI). This principle, applied throughout the \DPP\ architecture,
enables cross-lingual, cross-jurisdictional interoperability without requiring a
single natural language to dominate the system. The terminology is taken from
ontologies, which provide the semantic vocabulary for the \DPP\ knowledge
graphs.}
\label{fig:iri-multilingual}
\end{figure}

The terminology used in \DPP\ knowledge graphs is taken from ontologies, which
provide the semantic vocabulary for the \DPP\ data. The ontologies are
themselves identified by persistent \IRIs, ensuring that the meaning of \DPP\
data is unambiguously interpretable across the entire ecosystem. This layered
approach --- \IRIs\ for identification, ontologies for semantics, \SHACL\ for
structure --- follows the principle stated in the honorary lecture:
``\textit{everything must be ontologised}''~\cite{dpp4eu2026_day2}. The
ontologisation objective is an essentially global language that is adjusted to
local languages via mechanisms, while the user does not need to be confronted
with regulations. The user will need to know his product, but the system handles
the regulatory and semantic complexity transparently. This \IRI-based foundation
directly supports the external entity normalisation stage of the Ingestion
Governance Loop (Section~3.2.3): when a \DPP\ references a chemical substance, a
standard, or a certificate, the portal controller normalises the reference to a
canonical \IRI, ensuring consistent identification across all \DPPs\ in the
federation of national portals.

\subsection{Production as a Network: The Core Architectural Contrast}

A central insight of the proposed architecture is that production is a network,
not a linear chain. As stated in the honorary lecture at
\DPPforEU~2026~\cite{dpp4eu2026_day2}: ``\textit{production is a network. It's
at the very best case a tree. You have an awful lot of things that are being
assembled into something. But it's definitely not linear}''. This observation
has profound architectural consequences. The traditional approach of hardcoding
a static supply chain into a single product's \DPP\ --- whether as a flat-file
payload or as a blockchain record --- fundamentally misrepresents the structure
of production. A flat-file \DPP\ embeds all upstream data in a single document,
creating a monolithic record that is stale on creation and impossible to update
consistently. A blockchain \DPP\ replicates the supply chain as an append-only
ledger, creating unnecessary redundancy and immutability where dynamic
updatability is required. Both approaches treat the supply chain as a linear
chain; neither captures the reality that production is a network.

The Networked Individual \DPP\ model proposes a fundamentally different mapping:
the production network maps to a \textbf{\DPP\ graph network}, where each
product in the production network has its own \DPP, and the \DPPs\ are linked
via semantic properties (\texttt{dpp:hasPrecursor}, \texttt{dpp:derivedFrom}).
The supply chain information is not embedded in any single \DPP; it is
\textit{extracted from the \DPP\ repository by following the linked
\DPPs}~\cite{dpp4eu2026_day2}. Figure~\ref{fig:product-network} shows the
production network in terms of products and processes, and
Figure~\ref{fig:dpp-network} shows the corresponding \DPP\ graph network.

\begin{figure}[htbp]
\centering
\includegraphics[width=0.9\linewidth]{%
figures/2025-10-16__DigiPass_DPP_handling-process_network_flows.pdf}
\caption{The production network in terms of products and processes. Each process
(A--F) produces a product, and products flow between processes as inputs and
outputs. The network is a directed acyclic graph, not a linear chain: products
are assembled from multiple inputs, and a single product may be used in multiple
downstream processes. Each product flow has an associated \DPP.}
\label{fig:product-network}
\end{figure}

\begin{figure}[htbp]
\centering
\includegraphics[width=0.9\linewidth]{%
figures/2025-10-16__DigiPass_DPP_handling-process_network_DPP.pdf}
\caption{The corresponding \DPP\ graph network. Each product in the production
network has its own \DPP, stored as a knowledge graph behind a national portal
in the federation. The \DPPs\ are linked via \texttt{dpp:hasPrecursor} (assembly
links) and \texttt{dpp:derivedFrom} (transformation links). The supply chain
information is extracted by traversing the linked \DPPs\ within the federation
of national portals --- not by embedding it in a single monolithic record. The
portal controller maintains and updates the link network as the production
network evolves.}
\label{fig:dpp-network}
\end{figure}

The contrast between the three approaches is summarised in
Table~\ref{tab:approach-comparison}.

\begin{table}[p]
\centering
\caption{Comparison of supply chain representation approaches.}
\label{tab:approach-comparison}
\begin{tabular}{p{0.18\textwidth} p{0.27\textwidth} p{0.27\textwidth}
p{0.22\textwidth}}
\toprule
 & \textbf{Flat-File \DPP} & \textbf{Blockchain \DPP} & \textbf{Networked Graph
 \DPP\ (Proposed)} \\
\midrule
Supply chain model & Linear chain, embedded in single document & Linear chain,
replicated as append-only ledger & Network (\DAG), represented as linked graphs
\\
\addlinespace
Data structure & Single flat payload per product & Blockchain transaction per
product & \RDF\ named graph per product, linked via semantic properties \\
\addlinespace
Update propagation & None --- data stale on creation & Append-only --- old data
preserved, new data added & Instant --- single update to component \DPP\
propagates to all parents upon resolution \\
\addlinespace
Data duplication & Full duplication across all parent \DPPs & Full replication
across all nodes & Zero duplication --- links, not data \\
\addlinespace
Query capability & Limited to single \DPP\ payload & Limited to chain traversal
& Full \SPARQL\ queries across the entire \DPP\ network \\
\addlinespace
Access control & Binary --- full payload or nothing & Binary --- on-chain or
off-chain & Granular --- branch-level \ABAC\ with \ODRL\ policies \\
\addlinespace
\IP\ protection & None --- all data exposed in payload & None --- all data on
ledger & Branch-level filtering --- proprietary branches protected \\
\addlinespace
Scalability & Degrades with supply chain depth & Degrades with replication
overhead & Scales --- each \DPP\ is independently manageable \\
\bottomrule
\end{tabular}
\end{table}

The key architectural distinction is this: in the flat-file and blockchain
approaches, the supply chain is \textit{embedded} in the product's \DPP. In the
networked graph approach, the supply chain is \textit{emergent} from the links
between individual \DPPs. The supply chain does not exist as a stored data
structure; it is constructed on demand by traversing the linked \DPPs\ within
the portal's triple-store. This design also enables capabilities that neither
the flat-file nor the blockchain approach can support. A \SPARQL\ query can
recursively follow the \texttt{dpp:hasPrecursor} links, pull data attributes
from each child graph, and compile a real-time, mathematically sound Life Cycle
Assessment (\LCA)~\cite{dpp4eu2026_day2}.

The \DPP\ network also captures energy flows, waste streams, and utilities. As
the honorary lecture explained, each stream in the production process ---
including waste and utilities --- can be associated with ``\textit{a \DPP\ or
something of that kind}'', though it need not be a complete passport; what
matters is that it carries ``\textit{the relevant
information}''~\cite{dpp4eu2026_day2}. For some flows, such as electricity from
the grid, no \DPP\ exists; instead, the \LCA\ links to \textbf{secondary data}
from external lifecycle inventory databases to capture the impacts of energy
production~\cite{dpp4eu2026_day2}. The combination of linked \DPPs\ for material
streams and secondary data for utility flows enables comprehensive
system-boundary \LCA: because the portal stores both the mass and energy data
associated with each production step, a query engine can traverse the network
and aggregate impacts across the entire system boundary~\cite{dpp4eu2026_day2}.


% The key architectural distinction is this: in the flat-file and blockchain
% approaches, the supply chain is \textit{embedded} in the product's \DPP. In
% the networked graph approach, the supply chain is \textit{emergent} from the
% links between individual \DPPs. The supply chain does not exist as a stored
% data structure; it is constructed on demand by traversing the linked \DPPs
% within the portal's triple-store. This design also enables capabilities that
% neither the flat-file nor the blockchain approach can support. A \SPARQL query
% can recursively follow the \texttt{dpp:hasPrecursor} links, pull data
% attributes from each child graph, and compile a real-time, mathematically
% sound Life Cycle Assessment (\LCA)~\cite{dpp4eu2026_day2}. The \DPP network
% also captures energy flows, waste streams, and utilities --- each material
% stream in the production process, including waste and utilities, can have its
% own \DPP, enabling comprehensive system-boundary \LCA~\cite{dpp4eu2026_day2}.

\subsection{Resolving the Link Rot Problem}

Under the portal-backed model, because the individual graph of each component
was pushed \textit{behind the federated portal} at the time of manufacture, the
link remains entirely resolvable forever. The query engine follows the link
within the federation of national portals, retrieving the necessary data fields
smoothly even decades later. This is in stark contrast to the \REO-hosted
approach, where a component manufacturer going out of business would return a
\texttt{404 Not Found} error when a query engine tries to resolve
\texttt{dpp:hasPrecursor}. The parent product's carbon footprint calculation or
material circularity calculation would fail because the component data is simply
unreachable. The conference discussion on data persistency (see Section~2.1)
directly motivates this design. The portal-backed architecture answers this
question structurally rather than procedurally: data persists because it is
stored in a governed repository, not at an unmanaged endpoint that may disappear
at any time.

\subsection{Resolving the Data Update and Consistency Problem}

The networked model also resolves the data update problem inherent in monolithic
\DPPs. In the monolithic model, if a component's data is embedded in 100 parent
product \DPPs, updating the component data requires updating 100 separate \DPPs
--- a logistically impossible task across organisational boundaries. In the
networked model, the component data lives in exactly one place (the component's
own \DPP\ behind the portal), and all 100 parent products reference it by link.
A single update to the component \DPP\ propagates instantly to all parent
products upon resolution. This design also addresses the concern raised at the
conference about data correctness and dynamic updates: ``\textit{how we are
going to solve this issue or to guarantee the sustainability and working of}''
the data over time~\cite{dpp4eu2026_day2}. The architectural answer is that the
portal maintains a single authoritative copy of each \DPP, with versioning and
provenance metadata, and all references resolve to the current authoritative
version.
